\scanpage{35}
\mspage{16}
L'équation à satisfaire par $(\lambda,\eta)$ se réduit, on a vu, à $\unsure{\rho|\{}\ \widetilde\alpha\rho^{-1}(\widetilde\alpha)\rho^{-2}(\widetilde\alpha)$, avec $\widetilde\alpha$ donné en fonction de $(\lambda,\eta)$ par la première \deleted{relation} \inserted{formule} (55) --- mais cette relation \inserted{en $\widetilde\alpha$} est conséquence de la seconde relation (55), donnant $\widetilde\alpha$ en termes de $\xi$. Il reste donc, comme relation reliant $\lambda,\eta,\xi$, que la relation (55) sous la forme (où ne figure plus $\unsure{\widehat\alpha|\widetilde\alpha}$)
\[
(56)\qquad\eta^{-1}\varepsilon_0(\lambda\eta)=\xi\rho^{-1}(\xi)^{-1}
\]
Je vais la mettre sous une forme plus familière pour moi, en introduisant d'abord les quantités précédentes
\[
(57)\qquad\left\{\begin{aligned}
\beta&=\sigma_\infty(\eta)\eta^{-1}
&&\bigl(u(b_i^\omega)=b_i^\omega\varphi_{\unsure{i\omega|\omega i}}^\omega(\beta)\bigr)\\
\alpha&=\sigma_\infty(\lambda)\beta
&&\bigl(u(\widetilde b_i^\omega)=\widetilde b_i^\omega\varphi_{\omega i}^\omega(\alpha)\bigr)\\
\widetilde\alpha&=\eta^{-1}\varepsilon_0(\lambda\eta)=\xi\rho^{-1}(\xi)^{-1}
&&\bigl(u(\widetilde a_i^\omega)=\psi_{\omega i}^\omega(\widetilde\alpha)\widetilde a_i^\omega\bigr)
\end{aligned}\right.
\]
et en les modifiant ainsi, ainsi que $\xi$:
\[
(58)\qquad\left\{\begin{aligned}
\beta'&=\beta^{-1}\inserted{=\sigma_\infty(\beta)}=\eta\sigma_\infty(\eta)^{-1}\\
\alpha'&=\alpha^{-1}=\beta^{-1}\sigma_\infty(\lambda)^{-1}=\beta'\sigma_\infty^{-1}(\lambda^{-1})\\
\deleted{\xi'}&\deleted{=\unsure{\eta}\,\unsure{\eta|\eta^{-1}}\xi}
\end{aligned}\right.
\]
\margin{donc $\beta'=\alpha'\sigma_\infty(\lambda)$}
\deleted{Je vais \unsure{ainsi|aussi} \unsure{considérer}, en introduisant \unsure{un mot} $\unsure{\gamma}$ \unsure{en} $\eta$.}
\deleted{\unsure{(59)}}
\deleted{Je dis que la relation (56) équivaut à}
\deleted{\unsure{(59)}}
En résumé, \unsure{$u$ est décrit} par les quantités \figure{\begin{tikzpicture}\node[draw,ellipse] {$\lambda,\eta,\xi'$};\end{tikzpicture}} $\in\widehat\Pi_{0,3}^{\unsure{3}}$ (équivalt à la donnée de $\lambda,\eta,\xi$), auxquelles j'associe $\alpha',\beta'$ par
\[
(59)\qquad\begin{aligned}
\beta'&=\eta\sigma_\infty(\eta)^{-1}\deleted{\unsure{\sigma_0(\lambda\eta)}}\\
\deleted{\unsure{d'}}\alpha'&=\deleted{\unsure{\beta}}\beta'\sigma_\infty(\lambda^{-1})
\qquad\text{i.e.}\quad\beta'=\alpha'\sigma_\infty(\lambda)
\end{aligned}
\]
\margin{34}
